% Source: https://artofproblemsolving.com/community/c3h2279633_what_does_superdeficient_mean
% Stable identity: aops-topic-2279633-q01
\textbf{(a)} $f(2)=1+2=3$ and $f(3)=1+3=4=2+2$, so $2$ is superdeficient.

\textbf{(b)} For $n=1$, $f(f(1))=f(1)=1\ne3$, so $1$ is not superdeficient.

For $n\ge2$, the numbers $1$ and $n$ are different divisors of $n$, so $f(n)\ge n+1$, with equality exactly when $1$ and $n$ are the only divisors, that is, when $n$ is prime. Apply this twice; since $f(n)\ge n+1\ge3$,
\[f(f(n))\ge f(n)+1\ge(n+1)+1=n+2 .\]
So $f(f(n))=n+2$ forces equality in both steps: $f(n)=n+1$, so $n$ is prime, and $f(f(n))=f(n)+1$, so $f(n)=n+1$ is prime as well. Two consecutive integers $n$ and $n+1$ that are both prime must be $2$ and $3$, because one of them is even and $2$ is the only even prime ($n+1=2$ would give $n=1$, which is not prime). Hence $n=2$.

The only superdeficient positive integer is $\boxed{2}$.
